\section{MRV Design and Workflow}
\label{sec:design}

\subsection{Overview}

MRV processes each execution slice after the base protocol commits its vertices and the exporter produces the slice, but before those vertices are passed to the execution layer. For a slice $S$, MRV maintains creator-level visibility state while waiting for the exporter to expose the snapshot $P_{k^\star_S}$ defined in Section~\ref{sec:structural_evidence}. At that
snapshot, MRV fixes the visibility table for $S$, evaluates eligibility, compares causally incomparable vertices over their fixed pair windows, and builds $G_S$. Finally, it orders the SCCs of $G_S$, derives a deterministic total order within each SCC while preserving its causal edges, and outputs $\prec_S$. Figure~\ref{fig:mrv-overview} summarizes this workflow.

The vertices in $S$ form the ordering domain, whereas the committed prefix $P_{k^\star_S}$ forms the evidence domain. Descendants committed after $S$ was exported can appear in this prefix and provide evidence, but they never become members of $S$. MRV maintains state for all unsealed slices, but releases their execution orders in the slice order fixed by the exporter.

The protocol moves through three representations. For each round, a visibility profile records how many distinct creators have a canonical vertex whose ancestry contains the target. The pairwise comparator turns two such profiles into at most one ordering edge. The graph stage combines ordering edges with hard causal edges, computes the SCCs of the resulting graph, and derives a slice-wide order. This progression from per-vertex evidence to pairwise edges and then to a slice-wide order is described next.


\subsection{Snapshot Evidence Extraction}
\label{sec:evidence_extraction}

When $S$ is exported at event $o(S)$, MRV records its predetermined round horizon $T_r(S)=\max_{X\in S}r(X)+W$. It selects the first committed prefix at or after $o(S)$ whose highest round reaches this horizon. This prefix is exactly $P_{k^\star_S}$. The rule depends on the common prefix sequence, not on the physical time at which a replica observes a prefix.

For every target vertex $X\in S$, MRV obtains $C_X^{(k^\star_S)}(t)$ by examining the canonical vertex $Y^{(k^\star_S)}_{c,t}$ of each creator at round $t$. A creator contributes one unit if that vertex exists and has $X$ in its ancestry. The ancestry-closure property of the exporter makes this query computable from the selected prefix. A missing creator-round vertex contributes zero.
For brevity in the pseudocode, we write $\mathbf{C}_S[X,t]=C_X^{(k^\star_S)}(t)$ for every $X\in S$ and every round $t$ used by $X$'s eligibility interval or by a pair window involving $X$.

Each creator's contribution at a given round is binary. For example, if the canonical vertices of three creators at round $t$ contain $X$ in their ancestry, then $C_X^{(k^\star_S)}(t)=3$, regardless of how many paths from those vertices reach $X$. The sequence of these round-indexed counts is the visibility profile of $X$.

MRV evaluates $\mathsf{Elig}_S(X)$ over the fixed interval $\{r(X),\ldots,r(X)+W\}$. Eligibility checks whether the selected snapshot contains a round in which sufficiently many creators have $X$ in their vertices' ancestry. It does not by itself order $X$ against another vertex. A threshold crossing is only an eligibility witness and does not shorten any pair window. In the prototype, $q_{\mathrm{vis}}=2f+1$, so a witness includes canonical vertices from at least $f+1$ correct creators. All counts and eligibility flags for $S$ are evaluated on $P_{k^\star_S}$; later prefixes do not revise them.

Figure~\ref{fig:mrv-visibility} illustrates how creator-round observations in the fixed snapshot become visibility profiles and eligibility flags.


% TODO: redraw rq3-2.pdf for the fixed-window model. Remove stopping times and maturity, and mark the common snapshot horizon T_r(S).
%Creator-level visibility in the fixed snapshot. Each matrix cell records whether one creator's canonical vertex at that round has the target vertex in its ancestry. MRV sums at most one observation per creator and round to obtain the visibility profile, then evaluates eligibility over the target's fixed interval.
\begin{figure}[htbp]
    \centering
    \includegraphics[width=\columnwidth]{rq3-2.pdf}
    \caption{Structural visibility accumulation and AUF stopping times. Matrix cells indicate whether creator-round AUFs see $A$, $B$, both, or neither. MRV counts distinct creators that see each target AUF and fixes $h_X$ at the $2f+1$ maturity threshold or the observation cap $r(X)+W_{\max}$.}
    \label{fig:mrv-visibility}
\end{figure}


\subsection{Pairwise Comparison}
\label{sec:pairwise_comparison}

MRV does not compare vertices that are already ordered by ancestry. Their order is represented by hard causal edges during graph construction. For each causally incomparable pair $A,B\in S$, MRV uses the fixed window $\mathcal{W}(A,B)=\{s(A,B)+1,\ldots,s(A,B)+W\}$. Every value in this window is read from the same snapshot $P_{k^\star_S}$.

At a round $t$ in this window, a positive $\Delta_S(A,B,t)$ means that the visibility count of $A$ exceeds that of $B$; a negative value means the reverse. MRV does not sum these differences across rounds. It records whether either directional threshold is crossed at least once during the complete window.

Algorithm~\ref{alg:pair_edge} implements the predicate $\mathsf{Edge}_S$ from Section~\ref{sec:structural_property}. Both vertices must be eligible. MRV adds $A\to B$ only if the window contains a threshold crossing for $A$ over $B$ and no crossing in the reverse direction. If both directions cross, the result is \textsc{Conflict}; if neither crosses, it is \textsc{NoSignal}. If either vertex is ineligible, the result is \textsc{Ineligible}. These reason labels support analysis and measurement; all three cases add no ordering edge. The comparator is the only stage that creates ordering edges; later stages assemble these edges and complete the remaining choices.

Figure~\ref{fig:mrv-pairwise} shows the conversion from two visibility profiles to one of these four outcomes.


% TODO: redraw rq3-3.pdf for the fixed pair window. Remove maturity and the endogenous pair horizon, and show Ineligible, Conflict, and NoSignal.
% Pairwise comparison over a fixed window. A failed eligibility check yields \textsc{Ineligible}. For two eligible, causally incomparable vertices, MRV examines $\Delta_S(A,B,t)$ from $s(A,B)+1$ through $s(A,B)+W$. A threshold crossing in only one direction yields an ordering edge; crossings in both directions yield \textsc{Conflict}, and no crossing yields \textsc{NoSignal}.
\begin{figure}[htbp]
    \centering
    \includegraphics[width=.9\columnwidth]{rq3-3.pdf}
    \caption{Pairwise comparison and verdict extraction. For a mature pair $(A,B)$, MRV evaluates visibility deltas only after both AUFs coexist. An unopposed $f+1$ crossing yields $A\to B$; conflict or no strong signal yields abstention.}
    \label{fig:mrv-pairwise}
\end{figure}

\begin{algorithm}
\caption{Ordering edge for a causally incomparable pair}
\label{alg:pair_edge}
\begin{algorithmic}[1]
\Require $A,B\in S$; eligibility map $\mathsf{Elig}_S$;
         visibility table $\mathbf{C}_S$; threshold $\theta$
\If{$\neg\mathsf{Elig}_S(A)\lor\neg\mathsf{Elig}_S(B)$}
    \State \Return $\mathtt{NoEdge(Ineligible)}$
\EndIf
\State $pos_{AB}\gets\mathsf{false}$; $pos_{BA}\gets\mathsf{false}$
\For{each $t\in\mathcal{W}(A,B)$}
    \State $\delta\gets\mathbf{C}_S[A,t]-\mathbf{C}_S[B,t]$
    \State $pos_{AB}\gets pos_{AB}\lor(\delta\geq\theta)$
    \State $pos_{BA}\gets pos_{BA}\lor(\delta\leq-\theta)$
\EndFor
\If{$pos_{AB}\land\neg pos_{BA}$}
    \State \Return $\mathtt{Edge}(A\to B)$
\ElsIf{$pos_{BA}\land\neg pos_{AB}$}
    \State \Return $\mathtt{Edge}(B\to A)$
\ElsIf{$pos_{AB}\land pos_{BA}$}
    \State \Return $\mathtt{NoEdge(Conflict)}$
\Else
    \State \Return $\mathtt{NoEdge(NoSignal)}$
\EndIf
\end{algorithmic}
\end{algorithm}

The comparator records whether each directional threshold is crossed at least once; it does not accumulate visibility differences across rounds. A crossing is therefore a binary condition, and remaining above the threshold for several rounds adds no weight. The fixed endpoint $s(A,B)+W$ ensures that every eligible, causally incomparable pair is evaluated over the complete post-coexistence window; eligibility never shortens this window.

\subsection{Graph Construction and Linearization}
\label{sec:graph_linearization}

MRV constructs $G_S$ with vertex set $S$. For each unordered pair of distinct vertices $\{A,B\}\subseteq S$, MRV first checks whether the two vertices are causally comparable. If $B\in\mathsf{Anc}(A)$, it adds the hard causal edge $B\to A$; if $A\in\mathsf{Anc}(B)$, it adds $A\to B$. Otherwise, it runs Algorithm~\ref{alg:pair_edge} and adds the returned ordering edge, if any. The \textsc{Ineligible}, \textsc{Conflict}, and \textsc{NoSignal} outcomes add no edge.

Because ordering edges are derived independently for different pairs, the combined causal-and-ordering graph may contain cycles. MRV therefore computes the SCCs of $G_S$ and condenses them into a
DAG. Condensation replaces each SCC with one node and retains the edges between different SCCs. For example, the edges $A\to B$, $B\to C$, and $C\to A$ place all three vertices in one SCC because no total order can preserve all three edges. MRV provides no strict ordering guarantee for ordering edges inside the same SCC. Edges between different SCCs remain constraints in the condensation DAG.

MRV uses the fixed key $\kappa(X)=(r(X),\mathsf{cr}(X),\mathsf{dig}(X))$ to break ties at both stages of the linearization. For each SCC $C$, it defines $\kappa(C)=\min_{X\in C}\kappa(X)$. The condensation DAG is ordered by repeatedly selecting a zero-indegree SCC. When multiple SCCs have zero indegree, MRV selects the one with the smallest $\kappa(C)$.

Within each SCC, MRV orders vertices using only its hard causal subgraph. Although the full SCC may contain cycles, its hard causal subgraph is acyclic because the ancestry relation of the underlying DAG is acyclic. MRV repeatedly selects a zero-indegree vertex; when multiple vertices have zero indegree, it selects the one with the smallest $\kappa(X)$. It then concatenates the per-SCC vertex orders in the topological order of the SCCs. Consequently, the resulting $\prec_S$ preserves every hard causal edge and every ordering edge whose endpoints lie in different SCCs.

\begin{algorithm}
\caption{MRV ordering for one execution slice}
\label{alg:order_slice}
\begin{algorithmic}[1]
\Require Slice $S$, snapshot $P_{k^\star_S}$, $W$,
         $q_{\mathrm{vis}}$, $\theta$, and the fixed key $\kappa$
\State Read $\mathbf{C}_S$ for $P_{k^\star_S}$ from the maintained
       visibility state
\State Evaluate $\mathsf{Elig}_S(X)$ for every $X\in S$
\State Initialize $G_S$ with vertex set $S$ and no edges
\For{each unordered pair $\{A,B\}\subseteq S$}
    \If{$B\in\mathsf{Anc}(A)$}
        \State Add $B\to A$ to $G_S$
    \ElsIf{$A\in\mathsf{Anc}(B)$}
        \State Add $A\to B$ to $G_S$
    \Else
        \State $v\gets
          \Call{PairEdge}{A,B,\mathsf{Elig}_S,\mathbf{C}_S,\theta}$
        \If{$v=\mathtt{Edge}(U\to V)$}
            \State Add $U\to V$ to $G_S$
        \EndIf
    \EndIf
\EndFor
\State $\mathcal{C}\gets\Call{SCC}{G_S}$
\State $G_{\mathcal{C}}\gets\Call{Condense}{G_S,\mathcal{C}}$
\State $\langle C_1,\ldots,C_\ell\rangle\gets
       \Call{KeyedTopologicalSort}{G_{\mathcal{C}},\kappa(C)}$
\State $\prec_S\gets\langle\,\rangle$
\For{$j=1$ \textbf{to} $\ell$}
    \State $G^{\mathsf{causal}}_{C_j}\gets$
           the hard causal subgraph on $C_j$
    \State $L_j\gets
      \Call{KeyedTopologicalSort}
      {G^{\mathsf{causal}}_{C_j},\kappa(X)}$
    \State Append $L_j$ to $\prec_S$
\EndFor
\State \Return $\langle C_1,\ldots,C_\ell\rangle,\prec_S$
\end{algorithmic}
\end{algorithm}

\textsc{KeyedTopologicalSort} repeatedly selects a zero-indegree node, appends it to the output, and removes it together with its outgoing edges. When multiple zero-indegree nodes are available, it selects the one with the smallest supplied key. MRV uses $\kappa(C)$ for the condensation DAG and $\kappa(X)$ for the causal subgraph inside each SCC. This rule makes the SCC order and every per-SCC vertex order deterministic.

\subsection{Replica Procedure and Cost}
\label{sec:replica_procedure}

A replica registers a slice when it observes its export event. It records $o(S)$ and $T_r(S)$, then waits until the exporter exposes
$P_{k^\star_S}$. At that point, the replica runs Algorithm~\ref{alg:order_slice} once and seals both the SCC order and $\prec_S$. Later prefixes play no role in the output for that slice.


MRV creates visibility state for each vertex in $S_i$ when $S_i$ is exported and retains this state until $S_i$ is sealed. For each transition from $P_{k-1}$ to $P_k$, MRV processes each canonical vertex in $V_k\setminus V_{k-1}$ once, in nondecreasing round order. For each new vertex, it traverses the vertex's ancestry and updates the state of every tracked vertex it reaches. Because parent references point to lower rounds, a target is registered before any newly committed descendant can contribute to its visibility. When $S$ seals at $P_{k^\star_S}$, the cached entries for $S$ equal $\mathbf{C}_S$ evaluated on that snapshot. Sealing therefore does not require a full rescan of the committed prefix.

When $S_i$ seals, MRV places its SCC order and $\prec_{S_i}$ in an output queue and discards its slice-local visibility state, eligibility flags, pair outcomes, and graph. The sealed output remains in the queue until all earlier slices $S_1,\ldots,S_{i-1}$ have been released to execution. This preserves the exporter order even if local computation for a later slice finishes first.

MRV adds no protocol messages and relies on the base protocol's existing authentication checks. Let $m=|S|$ and let $L_S=T_r(S)-\min_{X\in S}r(X)+1$ be the round span relevant to $S$. For one
unsealed slice, the logical visibility table $\mathbf{C}_S$ contains $O(mL_S)$ values, and $G_S$ requires $O(m^2)$ space. Given constant-time access to $\mathbf{C}_S$, pairwise comparison takes $O(m^2W)$ time. The current prototype stores creator sets in an ordered round map, increasing the worst-case visibility state to $O(nmL_S)$ and direct pairwise-comparison time to $O(m^2W\log L_S)$.

Retaining all pair outcomes for analysis requires $O(m^2)$ additional entries. SCC computation and keyed topological sorting together take $O(|E(G_S)|+m\log m)$ time. The cost of visibility extraction depends on the committed-DAG ancestry index and the amount of ancestry data examined. These bounds exclude historical committed-DAG metadata.


\subsection{Security Scope}
\label{sec:security_scope}

With $\theta=f+1$, Byzantine creators alone cannot account for a threshold crossing at a given round because their total contribution to a visibility difference is at most $f$ in absolute value. A crossing therefore requires a positive net contribution of at least one unit from correct creators. Faulty replicas can choose their parents strategically, while message scheduling can affect which parents are available to correct replicas when they create their vertices. Both effects can change the ancestry recorded in the selected snapshot. The threshold therefore does not establish a majority preference among correct replicas, and the resulting ordering edge can still depend on adversarial scheduling.

The fixed key resolves only ties left by the topological constraints used at each stage. It makes these choices deterministic, but a key-based choice carries no guarantee about relative visibility or manipulation resistance.
MRV derives ordering edges only between vertices in the same slice; the exporter fixes the order between slices.




\subsection{Analysis and Guarantees}
\label{sec:analysis}

The following results establish the main properties of MRV under the exporter contract and definitions in Section~\ref{sec:model}. They cover the threshold interpretation, comparator correctness, deterministic slice ordering, ordering edges in the SCC order, causal refinement, sealing, rollback prevention, and transaction projection.

\begin{lemma}[Correct-Creator Contribution]
\label{lem:correct_contribution}
For $\theta=f+1$, if $\Delta_S(A,B,t)\geq\theta$, then
$\Delta_{\mathsf{cor}}(A,B,t)\geq1$.
\end{lemma}
\noindent\emph{Proof sketch.}
Write $\Delta_S(A,B,t)=\Delta_{\mathsf{cor}}(A,B,t) +\Delta_{\mathsf{byz}}(A,B,t)$.
At most $f$ Byzantine creators contribute at round $t$, and each contributes at most one unit in absolute value. Hence
$|\Delta_{\mathsf{byz}}(A,B,t)|\leq f$. Therefore,
$\Delta_{\mathsf{cor}}(A,B,t)
=\Delta_S(A,B,t)-\Delta_{\mathsf{byz}}(A,B,t)
\geq(f+1)-f=1$.

\begin{lemma}[Comparator Correctness]
\label{lem:comparator_correctness}
For causally incomparable $A,B\in S$, Algorithm~\ref{alg:pair_edge} returns $\mathtt{Edge}(A\to B)$ if and only if $\mathsf{Edge}_S(A,B)$ holds.
\end{lemma}
\noindent\emph{Proof sketch.}
The algorithm first checks $\mathsf{Elig}_S(A)\land\mathsf{Elig}_S(B)$.
It then evaluates every $t\in\mathcal{W}(A,B)$ using
$\mathbf{C}_S[A,t]-\mathbf{C}_S[B,t]
=\Delta_S(A,B,t)$.
Thus, $pos_{AB}$ is true exactly when $\mathsf{Pos}_S(A,B)$ holds. Moreover, $\mathcal{W}(A,B)=\mathcal{W}(B,A)$ and $\Delta_S(B,A,t)=-\Delta_S(A,B,t)$, so the test $\delta\leq-\theta$ is exactly the test for $\mathsf{Pos}_S(B,A)$. The algorithm returns $A\to B$ precisely when both vertices are eligible, $\mathsf{Pos}_S(A,B)$ holds, and $\mathsf{Pos}_S(B,A)$ does not, which is the definition of $\mathsf{Edge}_S(A,B)$.

\begin{theorem}[Deterministic Slice Ordering]
\label{thm:deterministic_slice_ordering}
Given the same committed-prefix sequence, execution-slice sequence, and protocol parameters, every slice $S$ for which $k^\star_S$ exists is sealed with the same ordered SCCs and the same total order $\prec_S$ at all correct replicas.
\end{theorem}
\noindent\emph{Proof sketch.}
The common execution-slice sequence fixes $S$ and its export event $o(S)$. The common committed-prefix sequence and protocol parameters then fix the same first index $k^\star_S$ and the same snapshot $P_{k^\star_S}$. The maintained visibility state at this snapshot equals the formally defined $\mathbf{C}_S$. Consequently, visibility counts, eligibility flags, pair outcomes, and $G_S$ are deterministic functions of the common inputs. The SCC partition of $G_S$ is unique. At each topological-sort step, the fixed total keys $\kappa(C)$ and $\kappa(X)$ select a unique zero-indegree candidate.
Hence all correct replicas obtain the same SCC order, the same per-SCC vertex orders, and the same concatenated order $\prec_S$.

\begin{theorem}[Ordering Edges in the SCC Order]
\label{thm:ordering_edges_scc_order}
For every ordering edge $A\to B$ in $G_S$ of a sealed slice $S$,
$\beta_S(A)\leq\beta_S(B)$. If $A$ and $B$ lie in different SCCs, then $A\prec_S B$.
\end{theorem}
\noindent\emph{Proof sketch.}
If $A$ and $B$ lie in different SCCs, the edge $A\to B$ becomes an edge from the SCC of $A$ to the SCC of $B$ in the condensation DAG. Every topological order of this DAG places the SCC of $A$ before the SCC of $B$. The concatenation of the per-SCC orders therefore places $A$ before $B$. If the endpoints lie in the same SCC, their SCC positions are equal, so $\beta_S(A)=\beta_S(B)$. MRV provides no strict ordering guarantee for that ordering edge.

\begin{theorem}[Global Causal Refinement]
\label{thm:global_causal_refinement}
For every $m$ such that $S_1,\ldots,S_m$ are sealed, define $\Pi_m=\prec_{S_1}\Vert\cdots\Vert\prec_{S_m}$. For all distinct
$A,B\in D_m$, if $B\in\mathsf{Anc}(A)$, then $B\prec_{\Pi_m}A$.
\end{theorem}
\noindent\emph{Proof sketch.}
Consider distinct $A,B\in D_m$ with $B\in\mathsf{Anc}(A)$. If they belong to the same slice $S$, then $G_S$ contains the hard causal edge $B\to A$. If its endpoints lie in different SCCs, the condensation-DAG order places the SCC of $B$ before the SCC of $A$. If they lie in the same SCC, the topological sort of that SCC's hard causal subgraph places $B$ before $A$. Hence $B\prec_S A$ in both cases. If $A$ and $B$ belong to different slices, the
exporter's causal-delivery property gives $\sigma(B)\leq\sigma(A)$. The slice indices are different, so
$\sigma(B)<\sigma(A)$, and the concatenation of the slice orders places $B$ before $A$.

\begin{theorem}[Conditional Sealing]
\label{thm:conditional_sealing}
If the exporter eventually exposes a prefix $P_k$ with $k\geq o(S)$ and $\rho(k)\geq T_r(S)$, then $k^\star_S$ exists and MRV seals $S$ after processing $P_{k^\star_S}$ and completing finite local computation.
\end{theorem}
\noindent\emph{Proof sketch.}
The stated condition makes the set defining $k^\star_S$ nonempty, so its least element exists. For every $X\in S$, the eligibility interval ends at $r(X)+W\leq T_r(S)$. For every $A,B\in S$, the pair window ends at $H(A,B)=\max(r(A),r(B))+W\leq T_r(S)$. Once
$P_{k^\star_S}$ has been processed, MRV reads the maintained visibility state for this snapshot and evaluates eligibility, pair outcomes, SCCs, and the keyed topological orders over a finite slice, graph, and set of rounds. The remaining local computation is therefore finite.

\begin{corollary}[No Rollback]
\label{cor:no_rollback}
Once MRV seals $S$, no later committed prefix changes its ordered SCCs or $\prec_S$.
\end{corollary}
\noindent\emph{Proof sketch.}
Both outputs are deterministic functions of the fixed snapshot
$P_{k^\star_S}$. At sealing, MRV fixes these outputs and discards the slice-local evidence and graph state. Later committed prefixes therefore do not recompute or revise the output of $S$.

\begin{corollary}[Conditional Transaction Projection]
\label{cor:transaction_projection}
Suppose transactions $tx_A$ and $tx_B$ each occur in exactly one vertex of $V^{\mathsf{exec}}$, namely distinct vertices $A$ and $B$ of the same sealed slice $S$. Suppose further that the execution layer processes complete vertices in MRV order. If $\beta_S(A)<\beta_S(B)$, then $tx_A$ executes before $tx_B$. In particular, this holds when an ordering edge $A\to B$ has endpoints in different SCCs.
\end{corollary}
\noindent\emph{Proof sketch.}
The strict inequality $\beta_S(A)<\beta_S(B)$ places the SCC of $A$ before the SCC of $B$. Because $\prec_S$ concatenates the per-SCC vertex orders, it places every vertex in the SCC of $A$ before every vertex in the SCC of $B$. The execution layer processes complete vertices in this order, so $tx_A$ executes before $tx_B$. An ordering edge whose endpoints lie in different SCCs gives the required strict SCC order by Theorem~\ref{thm:ordering_edges_scc_order}. An ordering edge inside one SCC alone provides no such transaction-order guarantee.